\section{Methodology: Risk-Based Segmentation}
\label{sec:method}

Our methodology brings risk suitability inside the recommender. The core idea is \emph{risk-based segmentation}: instead of a single global model, we route each customer to a sub-model trained only on customers of their risk band. Figure~\ref{fig:overview} gives a full picture of our methodology.

\begin{figure*}[t]
\centering
\resizebox{\textwidth}{!}{%
\begin{tikzpicture}[
  >=Stealth, font=\scriptsize,
  block/.style={draw, rounded corners=3pt, align=center, inner sep=5pt},
  inp/.style={draw, rounded corners=3pt, align=center, font=\scriptsize, inner sep=4pt, fill=gray!4},
  prev/.style={draw, thick, rounded corners=3pt, align=center, fill=red!8, text width=3.6cm, minimum height=1.6cm, inner sep=5pt, font=\scriptsize},
  sub/.style={draw, thick, rounded corners=2pt, align=center, font=\scriptsize, minimum width=3.3cm, minimum height=0.5cm, fill=red!8},
  evalblock/.style={draw, rounded corners=3pt, align=center, fill=blue!6, text width=4.0cm, inner sep=7pt},
  alt/.style={draw, dashed, rounded corners=3pt, align=center, font={\scriptsize\hyphenpenalty=10000}, fill=gray!12, inner sep=4pt, text width=3.6cm},
  io/.style={align=center, font=\scriptsize},
  grp/.style={draw, dashed, rounded corners=6pt, inner sep=9pt},
  arr/.style={->, thick},
  darr/.style={->, thick, dashed},
]

% --- layout anchor (top-left reference) ---
\coordinate (anchor) at (0,0);

% --- (b) Ours: sub-models (shifted right) ---
\node[sub, anchor=north west] (s0) at ([xshift=7.0cm, yshift=-4.3cm]anchor) {Conservative risk-band sub-model};
\node[sub, below=3pt of s0] (s1) {Income risk-band sub-model};
\node[sub, below=3pt of s1] (s2) {Balanced risk-band sub-model};
\node[sub, below=3pt of s2] (s3) {Aggressive risk-band sub-model};

% router centered on the stack
\node[block, anchor=east] (router) at ([xshift=-1.5cm]$(s0.west)!0.5!(s3.west)$) {Risk-band\\router};

% --- (a) Previous lane (kept near the top, decoupled from the stack) ---
\coordinate (prevpos) at (s0.north west |- anchor);
\node[prev, anchor=north west] (prev) at ([yshift=-1.2cm]prevpos) {\textbf{(a) Baseline}\\[2pt] A single risk band-agnostic \mbox{LightGCN}};

% --- shared customers input, boxed, anchored at the left edge ---
\coordinate (lanemid) at ($(prev.west)!0.5!(router.west)$);
\node[inp, anchor=west] (cust) at ([xshift=0.15cm]lanemid -| anchor) {\textbf{FAR-Trans Data}};

% --- levers ---
\node[alt, anchor=south] (alt2) at ([yshift=0.8cm]router.north) {\textbf{Approach 2 (data lever):} Regroup customers onto revealed risk-band};
\node[alt, anchor=north] (alt1) at ([yshift=-0.5cm]s3.south) {\textbf{Approach 1 (model lever):} With additional margin loss};

% group box for (b)
\begin{scope}[on background layer]
\node[grp, fit=(router)(alt2)(s0)(s3)(alt1), label={[font=\scriptsize\bfseries]below:(b) Methodology: Risk-based segmentation}] (seg) {};
\end{scope}

% --- evaluation terminal (right) ---
\node[evalblock, anchor=west] (eval) at ([xshift=2.2cm]seg.east |- cust) {\textbf{Evaluate on three axes:}\\[3pt] Return (ROI)\\ Accuracy (nDCG)\\ Risk suitability (PC)\\[4pt] With the harmonic mean of the three.};

% --- intra-(b) wiring ---
\draw[arr] (cust) |- (prev.west);
\draw[arr] (cust) |- (router.west);

\coordinate (fo) at ([xshift=0.6cm]router.east);
\draw[thick] (router.east) -- (fo);
\draw[arr] (fo) |- (s0.west);
\draw[arr] (fo) |- (s1.west);
\draw[arr] (fo) |- (s2.west);
\draw[arr] (fo) |- (s3.west);

\coordinate (bus) at ([xshift=0.8cm]$(s0.east)!0.5!(s3.east)$);
\draw[thick] (s0.east) -- (s0.east -| bus);
\draw[thick] (s1.east) -- (s1.east -| bus);
\draw[thick] (s2.east) -- (s2.east -| bus);
\draw[thick] (s3.east) -- (s3.east -| bus);
\draw[thick] (s0.east -| bus) -- (s3.east -| bus);

% --- both lanes converge to top-k, then evaluated ---
\coordinate (jct) at ([xshift=-1.1cm]eval.west);
\draw[thick] (prev.east) -| (jct);
\draw[thick] (bus) -| (jct);
\draw[arr] (jct) -- (eval.west) node[io, midway, above] {top-$k$};

% --- levers point at where they act ---
\draw[darr] (alt2) -- (router);
\draw[darr] (alt1.north) -- (s3.south);

\end{tikzpicture}%
}
\caption{Overview of risk-based segmentation. The band-agnostic baseline (a) and our segmentation (b) are built from the same FAR-Trans data and scored on the same three axes. In (b) a band router sends each customer to the LightGCN sub-model of their band, and two exclusive levers raise risk suitability: a data-side regrouping onto the revealed band or a model-side coherence margin.}
\label{fig:overview}
\end{figure*}

\subsection{Suitability Gap of Single Global Model}
The FAR-Trans LightGCN baseline is a single collaborative filter over one user-item bipartite graph that connects every customer to the assets they bought. It learns an embedding per customer and per asset, propagates them over the graph, and scores an asset for a customer by the inner product of their embeddings, training with the Bayesian Personalised Ranking (BPR) objective~\cite{rendle2012bprbayesianpersonalizedranking}. The model never sees the customer's risk band, so suitability is whatever the purchase graph happens to imply. Section~\ref{sec:analysis} showed that purchases are frequently discordant, so a band-agnostic model trained to reproduce them inherits that discordance.

\subsection{Risk-Based Segmentation}
We choose LightGCN as the backbone because its differentiable loss and per-customer embeddings let us condition on the risk band and, later, attach a coherence margin directly to the BPR objective. Prior work shows LightGCN readily carries domain-specific structure beyond plain BPR~\cite{Ghiye_2024}. Random Forest, the other FAR-Trans anchor, ranks assets from technical indicators alone and carries no customer-level signal to segment on, so we do not extend it.

Risk-based segmentation partitions the customer base by declared MiFID II band and trains one LightGCN sub-model per band: Conservative, Income, Balanced and Aggressive. Each sub-model's training graph contains only the interactions of customers in that band, so its embeddings specialise to that band's preferences. A band router serves each customer through the sub-model of their declared band. Customers with no declared band do not enter any sub-model at training and fall back to the Balanced sub-model at inference. Each sub-model is trained with the same BPR objective and edge-dropout regularisation as the baseline, changing only the population it learns from.

\subsection{Two Approaches for Raising Risk Suitability}
Segmentation specialises each sub-model to a band but still trains purely to reproduce purchases, which Section~\ref{sec:analysis} showed are themselves often discordant. We therefore study two exclusive levers for pushing risk suitability higher, one acting on the model and one on the data.

\subsubsection{Approach 1 (model lever): Steering the model via margin loss}
BPR trains the model to score a bought asset (the positive) above a random unobserved asset (the negative). Independently of this sampling, we add a coherence margin that draws, from the sub-model's asset universe, one coherent asset $i_c$ ($d(u, i_c) \leq 1$) and one discordant asset $i_d$ ($d(u, i_d) > 1$), and pushes the coherent score above the discordant one:
\begin{equation}
\mathcal{L}_{\text{coh}}(u) = \text{softplus}\!\left(s(u, i_d) - s(u, i_c)\right),
\end{equation}
where $s(u, i)$ is the inner-product score and $\text{softplus}(x) = \ln(1 + e^x)$ smoothly approximates the hinge $\max(0, x)$. The total objective adds this term to BPR with weight $\lambda \geq 0$:
\begin{equation}
\mathcal{L} = \mathcal{L}_{\text{BPR}} + \omega \|\Theta\|_2^2 + \lambda \cdot \mathcal{L}_{\text{coh}},
\end{equation}
where $\omega$ is the L2 weight decay on embedding parameters $\Theta$. BPR remains the primary ranking signal: setting $\lambda = 0$ recovers plain segmentation, while $\lambda > 0$ steers scores towards coherent assets. We use softplus rather than a hard hinge because it is differentiable everywhere and is the same form BPR itself uses, keeping the two objectives consistent.

\subsubsection{Approach 2 (data lever): Fixing the data via regrouping}
The second approach corrects the label at its source: before training, each banded customer is reassigned to the band implied by their purchases, taken as the lower median of the risk classes of the assets they bought (even-count ties resolve to the more conservative band). Segmentation then routes and trains on these regrouped bands, making the data coherent in the first place rather than penalising incoherence during training.